\documentclass[11pt]{article}

\usepackage[margin=1in]{geometry}
\usepackage[T1]{fontenc}
\usepackage[utf8]{inputenc}
\usepackage{lmodern}
\usepackage{amsmath,amssymb,mathtools}
\usepackage{enumitem}
\usepackage{array,booktabs,longtable,tabularx}
\usepackage{xcolor}
\usepackage{hyperref}
\usepackage{listings}
\usepackage{tcolorbox}
\tcbuselibrary{skins,breakable}
\usepackage{titlesec}

\hypersetup{
    colorlinks=true,
    linkcolor=blue!60!black,
    urlcolor=blue!60!black
}

\definecolor{myblue}{RGB}{20,55,90}
\definecolor{mygray}{RGB}{245,245,245}
\definecolor{mygreen}{RGB}{20,110,70}

\titleformat{\section}{\large\bfseries\color{myblue}}{\thesection.}{0.5em}{}
\titleformat{\subsection}{\normalsize\bfseries\color{myblue}}{\thesubsection.}{0.5em}{}

\lstdefinestyle{mystyle}{
    basicstyle=\ttfamily\small,
    backgroundcolor=\color{mygray},
    frame=single,
    breaklines=true,
    showstringspaces=false,
    columns=fullflexible
}

\lstset{style=mystyle}

\newtcolorbox{goalbox}{
    colback=blue!3,
    colframe=myblue,
    boxrule=0.8pt,
    arc=2pt
}

\newtcolorbox{warningbox}{
    colback=red!2,
    colframe=red!55!black,
    boxrule=0.8pt,
    arc=2pt
}

\newtcolorbox{intuitbox}{
    colback=green!3,
    colframe=mygreen,
    boxrule=0.8pt,
    arc=2pt
}

\begin{document}

\begin{center}
    {\LARGE \textbf{Detailed Execution Plan: AI-Assisted Macroeconomic Modeling Dashboard}}\\[0.5em]
    {\large ECO 317 -- Assignment \#2}\\[0.5em]
    \rule{0.9\linewidth}{0.5pt}
\end{center}

\section{Detailed Step-by-Step Plan}

\subsection{Step 1: Read the course notes and lock in the exact equations}
\textbf{Goal:} Fix the economics before any solver code exists.

\textbf{Do this:}
\begin{itemize}
    \item Make one shared ``model lock-in sheet'' with a separate block for Models 1, 2, and 3.
    \item For each model, write down:
    \begin{itemize}
        \item Bellman equation,
        \item budget or resource constraint,
        \item state variables,
        \item control variables,
        \item shock states,
        \item transition matrix,
        \item timing convention,
        \item variables that must be simulated and plotted.
    \end{itemize}
    \item Decide coding names now, for example \verb|a|, \verb|a_next|, \verb|k|, \verb|k_next|, \verb|z_idx|, \verb|P|.
    \item Add one line saying what each solver will return.
    \item Mark any place where the notes differ from a textbook version.
\end{itemize}

\textbf{Example workflow:}
\begin{itemize}
    \item Model 1 block:
    \begin{itemize}
        \item state = assets and income shock,
        \item control = next-period assets,
        \item consumption comes from the budget constraint.
    \end{itemize}
    \item Model 2 block:
    \begin{itemize}
        \item explicitly write whether output uses current capital or next-period capital.
    \end{itemize}
    \item Model 3 block:
    \begin{itemize}
        \item explicitly write whether assets belong in the lecture version.
    \end{itemize}
\end{itemize}

\textbf{Deliverable:}
\begin{itemize}
    \item One model lock-in sheet used by the entire team.
\end{itemize}

\textbf{Quick check:}
\begin{itemize}
    \item If someone asks ``what is the state, control, constraint, and timing?'' for any model, the answer should be immediate.
\end{itemize}

\subsection{Step 2: Set up the repository and development environment}
\textbf{Goal:} Make the project launch the same way on every laptop.

\textbf{Do this:}
\begin{itemize}
    \item Create the repository and folder tree.
    \item Add \verb|requirements.txt| first.
    \item Create a virtual environment.
    \item Install packages.
    \item Make a tiny \verb|app.py| test page to confirm Streamlit launches.
    \item Check that all teammates can run the app locally.
\end{itemize}

\textbf{Files to create immediately:}
\begin{itemize}
    \item \verb|app.py|
    \item \verb|requirements.txt|
    \item \verb|README.md|
    \item \verb|config.py|
    \item \verb|assets/style.css|
    \item \verb|models/|, \verb|solvers/|, \verb|simulation/|, \verb|utils/|, \verb|tests/|
\end{itemize}

\textbf{Example workflow:}
\begin{itemize}
    \item Use a minimal smoke test first:
\end{itemize}
\begin{lstlisting}
import streamlit as st

st.title("ECO 317 Project 2")
st.write("Environment test passed.")
\end{lstlisting}
\begin{itemize}
    \item Run \verb|streamlit run app.py| before writing any model logic.
    \item Fix missing-package or Python-version problems now rather than during model development.
\end{itemize}

\textbf{Deliverable:}
\begin{itemize}
    \item A clean repo that installs with one command and launches with one command.
\end{itemize}

\textbf{Quick check:}
\begin{itemize}
    \item \verb|pip install -r requirements.txt| works.
    \item \verb|streamlit run app.py| opens in a browser without errors.
\end{itemize}

\subsection{Step 3: Build the project skeleton before filling in model logic}
\textbf{Goal:} Make the architecture visible before the hard coding begins.

\textbf{Do this:}
\begin{itemize}
    \item Create empty modules with docstrings and placeholder functions.
    \item Define consistent function signatures.
    \item Decide one solver return structure for all models.
    \item Decide where parameter dictionaries or dataclasses will live.
\end{itemize}

\textbf{Example workflow:}
\begin{itemize}
    \item Put placeholders like this into each model file:
\end{itemize}
\begin{lstlisting}
def solve_consumption_savings(params):
    """Return value function, policies, grids, and diagnostics."""
    raise NotImplementedError
\end{lstlisting}
\begin{itemize}
    \item Decide early that each solver returns keys such as:
    \begin{itemize}
        \item \verb|value_function|,
        \item \verb|policy_indices|,
        \item \verb|policy_levels|,
        \item \verb|grids|,
        \item \verb|diagnostics|.
    \end{itemize}
    \item This prevents each teammate from inventing a different output format.
\end{itemize}

\textbf{Deliverable:}
\begin{itemize}
    \item A modular skeleton where every major task already has a home.
\end{itemize}

\textbf{Quick check:}
\begin{itemize}
    \item A new teammate should be able to glance at the file tree and understand where each piece belongs.
\end{itemize}

\subsection{Step 4: Build shared utilities first}
\textbf{Goal:} Build the reusable numerical pieces once.

\textbf{Do this:}
\begin{itemize}
    \item In \verb|utils/grids.py|, write functions to create linear grids and optionally curved grids.
    \item Add clipping and bounds-check helpers.
    \item In \verb|utils/markov.py|, write transition-matrix validation and shock-path simulation.
    \item In \verb|utils/interpolation.py|, wrap \verb|numpy.interp| so simulation and forecasting can evaluate policies off-grid.
    \item In \verb|simulation/moments.py|, write functions for mean, variance, lag-1 autocorrelation, and correlations.
\end{itemize}

\textbf{Example workflow:}
\begin{itemize}
    \item Start the asset or capital grid with something like:
\end{itemize}
\begin{lstlisting}
grid = np.linspace(x_min, x_max, n_points)
\end{lstlisting}
\begin{itemize}
    \item For the Markov chain, create a validator that checks:
    \begin{itemize}
        \item correct shape,
        \item probabilities are nonnegative,
        \item each row sums to 1.
    \end{itemize}
    \item For interpolation, create one helper so all models use the same policy-evaluation rule when the state is off-grid.
\end{itemize}

\textbf{Important detail:}
\begin{itemize}
    \item The transition matrix should be validated carefully: rows must sum to 1, probabilities must be nonnegative, and the shock-state ordering must be consistent everywhere.
\end{itemize}

\textbf{Deliverable:}
\begin{itemize}
    \item A reusable numerical toolkit used by all three models.
\end{itemize}

\textbf{Quick check:}
\begin{itemize}
    \item You can generate a valid grid, simulate a two-state Markov chain, and compute moments on a fake time series.
\end{itemize}

\subsection{Step 5: Implement the generic VFI engine}
\textbf{Goal:} Create one Bellman-iteration backbone that all models can use.

\textbf{Do this:}
\begin{itemize}
    \item Write the main VFI loop in \verb|solvers/vfi.py|.
    \item Allow the model file to supply model-specific pieces such as utility, feasibility, state transition logic, and candidate controls.
    \item Track convergence distance at each iteration.
    \item Return diagnostics such as number of iterations, final error, and whether convergence was achieved.
    \item Store policy indices so simulation does not need to re-solve the Bellman problem later.
\end{itemize}

\textbf{Example workflow:}
\begin{itemize}
    \item Start with:
\end{itemize}
\begin{lstlisting}
V = np.zeros((n_grid, n_shocks))
for it in range(max_iter):
    V_new = np.empty_like(V)
    policy_idx = np.empty_like(V, dtype=int)
\end{lstlisting}
\begin{itemize}
    \item For each state:
    \begin{itemize}
        \item generate candidate controls,
        \item remove infeasible choices,
        \item compute current utility,
        \item compute expected continuation value,
        \item maximize over controls,
        \item store the best index.
    \end{itemize}
    \item At the end of each iteration compute:
\end{itemize}
\begin{lstlisting}
err = np.max(np.abs(V_new - V))
\end{lstlisting}
\begin{itemize}
    \item Stop when \verb|err < tol|.
\end{itemize}

\textbf{Important implementation choice:}
\begin{itemize}
    \item Keep the generic solver reusable, but not so abstract that debugging becomes impossible.
\end{itemize}

\textbf{Deliverable:}
\begin{itemize}
    \item A reusable solver backbone that every model can call.
\end{itemize}

\textbf{Quick check:}
\begin{itemize}
    \item The solver converges on a tiny toy problem with known behavior.
\end{itemize}

\subsection{Step 6: Add hard feasibility checks before any real model runs}
\textbf{Goal:} Ensure impossible choices never enter the maximization.

\textbf{Do this:}
\begin{itemize}
    \item For each model, define a feasibility function.
    \item Any infeasible choice should be excluded from the maximization.
    \item Never let negative consumption pass silently.
    \item Decide whether infeasible values get \verb|-np.inf| utility or are filtered out before maximization.
    \item Add assertions in testing so infeasible policies are caught immediately.
\end{itemize}

\textbf{Example workflow:}
\begin{itemize}
    \item In Model 1, compute:
\end{itemize}
\begin{lstlisting}
c = (1 + r) * a + y - a_prime
feasible = c > 0
\end{lstlisting}
\begin{itemize}
    \item In the labor model also require:
    \begin{itemize}
        \item \(0 \leq L \leq 1\),
        \item \(1-L > 0\) when utility needs positive leisure,
        \item next-period states stay within bounds.
    \end{itemize}
    \item Either drop infeasible controls or assign them \verb|-np.inf|.
\end{itemize}

\textbf{Deliverable:}
\begin{itemize}
    \item Reliable numerical safeguards.
\end{itemize}

\textbf{Quick check:}
\begin{itemize}
    \item With extreme parameter values, the solver should still fail cleanly or exclude bad choices, not produce nonsense.
\end{itemize}

\subsection{Step 7: Implement Model 1 completely before touching Models 2 and 3}
\textbf{Goal:} Use Model 1 as the full prototype pipeline.

\textbf{Do this:}
\begin{itemize}
    \item Create the asset grid.
    \item Define low and high income states.
    \item Define the \(2 \times 2\) Markov transition matrix.
    \item Implement the utility function, including the log case when \verb|sigma = 1|.
    \item For each current asset and current income state, evaluate all feasible choices of next-period assets.
    \item Compute the expected continuation value using the transition matrix.
    \item Solve for the value function, savings policy, and implied consumption policy.
\end{itemize}

\textbf{Example workflow:}
\begin{itemize}
    \item Start with a grid like:
\end{itemize}
\begin{lstlisting}
a_min, a_max, n_a = 0.0, 20.0, 200
asset_grid = np.linspace(a_min, a_max, n_a)
\end{lstlisting}
\begin{itemize}
    \item Define income states such as \verb|y_vals = [y_low, y_high]|.
    \item Define the transition matrix:
\end{itemize}
\begin{lstlisting}
P = np.array([[0.9, 0.1],
              [0.1, 0.9]])
\end{lstlisting}
\begin{itemize}
    \item For each \((a, y)\):
    \begin{itemize}
        \item loop through candidate \(a'\) values,
        \item compute consumption,
        \item remove infeasible choices,
        \item evaluate Bellman value,
        \item save best \(a'\) and implied \(c\).
    \end{itemize}
\end{itemize}

\textbf{Economic checks:}
\begin{itemize}
    \item Consumption should be positive everywhere.
    \item Savings should typically rise with assets.
    \item Higher patience should generally increase saving.
    \item Higher risk aversion should generally smooth consumption more.
\end{itemize}

\textbf{Deliverable:}
\begin{itemize}
    \item One fully solved model with believable policies.
\end{itemize}

\textbf{Quick check:}
\begin{itemize}
    \item You can plot \(a'(a,y)\) and \(c(a,y)\) for both income states and the curves look smooth and economically sensible.
\end{itemize}

\subsection{Step 8: Build simulation and moments for Model 1}
\textbf{Goal:} Turn the solved policies into a usable time series.

\textbf{Do this:}
\begin{itemize}
    \item Choose an initial asset level and initial shock state.
    \item Simulate a shock path using the Markov chain.
    \item Use the saved policy functions each period.
    \item Interpolate when the current state falls between grid points.
    \item Store the full time paths for assets, consumption, and income.
    \item Convert the arrays into a DataFrame.
    \item Compute mean, variance, lag-1 autocorrelation, and correlation with income.
\end{itemize}

\textbf{Example workflow:}
\begin{itemize}
    \item Fix:
    \begin{itemize}
        \item \(T = 100\) or \(200\),
        \item \(a_0\),
        \item initial shock state,
        \item random seed.
    \end{itemize}
    \item At each period:
    \begin{itemize}
        \item read current state,
        \item interpolate the policy if the asset level is off-grid,
        \item compute next assets,
        \item recover consumption from the budget constraint,
        \item store assets, consumption, and income.
    \end{itemize}
    \item Build a DataFrame with columns like \verb|assets|, \verb|consumption|, \verb|income|.
\end{itemize}

\textbf{Important detail:}
\begin{itemize}
    \item Do not recompute the maximization during simulation. Simulation should only read the solved policy functions.
\end{itemize}

\textbf{Deliverable:}
\begin{itemize}
    \item A complete Model 1 pipeline from solver to time series to moments.
\end{itemize}

\textbf{Quick check:}
\begin{itemize}
    \item With the same random seed, the simulation should be reproducible.
\end{itemize}

\subsection{Step 9: Build forecast logic for Model 1 using the solved policy functions}
\textbf{Goal:} Forecast with policy rules instead of trend fitting.

\textbf{Do this:}
\begin{itemize}
    \item Write a forecast function that accepts the current state and a user-chosen future shock path.
    \item Use the actual policy rules to update the state period by period.
    \item Generate forecast paths for consumption, savings or assets, and income.
    \item Make the horizon configurable.
    \item Keep the forecast logic separate from stochastic simulation, since here the future shock path is user-controlled.
\end{itemize}

\textbf{Example workflow:}
\begin{itemize}
    \item Suppose the current state is ``assets = 4.2, shock = low income.''
    \item Suppose the chosen future shocks are \verb|[high, high, low, high]|.
    \item For each forecast step:
    \begin{itemize}
        \item evaluate the saved policy,
        \item update assets,
        \item move to the next user-chosen shock,
        \item store the forecasted variables.
    \end{itemize}
    \item Keep the forecast arrays separate from the stochastic simulation arrays.
\end{itemize}

\textbf{Deliverable:}
\begin{itemize}
    \item Forecast series that are genuinely implied by the solved model.
\end{itemize}

\textbf{Quick check:}
\begin{itemize}
    \item If the user chooses a sequence of high shocks, the forecast should reflect that path in a way consistent with the model's policy functions.
\end{itemize}

\subsection{Step 10: Implement Model 2 only after a timing checkpoint}
\textbf{Goal:} Build the Robinson Crusoe model without a timing error.

\textbf{Do this:}
\begin{itemize}
    \item Re-open the notes and verify the exact capital-output timing.
    \item Define the capital grid and TFP shock states.
    \item Implement the transition matrix for TFP.
    \item Write production and depreciation logic.
    \item Solve for the value function, next-period capital policy, and implied consumption policy.
    \item Add implied output and investment calculations for simulation and plots.
\end{itemize}

\textbf{Example workflow:}
\begin{itemize}
    \item Write a comment at the top of the file stating the timing convention taken from the notes.
    \item Start with:
\end{itemize}
\begin{lstlisting}
k_grid = np.linspace(k_min, k_max, n_k)
z_vals = np.array([z_low, z_high])
\end{lstlisting}
\begin{itemize}
    \item For each \((k, z)\):
    \begin{itemize}
        \item loop over candidate \(k'\),
        \item compute output,
        \item compute consumption,
        \item exclude infeasible choices,
        \item maximize Bellman value.
    \end{itemize}
    \item Save derived quantities such as output and investment for simulation plots.
\end{itemize}

\textbf{Economic checks:}
\begin{itemize}
    \item Output should be higher in the good TFP state.
    \item Capital policy should usually respond positively to better productivity.
    \item Consumption and investment should behave sensibly when TFP changes.
\end{itemize}

\textbf{Deliverable:}
\begin{itemize}
    \item A solved stochastic growth-style model with capital accumulation.
\end{itemize}

\textbf{Quick check:}
\begin{itemize}
    \item The model should produce reasonable paths for capital, output, consumption, and investment.
\end{itemize}

\subsection{Step 11: Implement Model 3 with labor-leisure trade-offs and confirm whether assets belong}
\textbf{Goal:} Build the labor model exactly as the lecture version intends.

\textbf{Do this:}
\begin{itemize}
    \item Confirm from the notes whether the state includes assets or just the wage shock and labor choice.
    \item Define wage states and the wage transition matrix.
    \item Define the labor grid on \([0,1]\).
    \item If assets are included, jointly search over labor and next-period assets.
    \item If assets are not included, solve the simpler static or dynamic labor-choice version from the notes.
    \item Recover labor, consumption, and savings policies if relevant.
    \item Simulate labor, wage, and consumption paths.
    \item Build at least one static intuition graph, such as labor against wage or labor against the preference parameter.
\end{itemize}

\textbf{Example workflow:}
\begin{itemize}
    \item Start with a labor grid:
\end{itemize}
\begin{lstlisting}
labor_grid = np.linspace(0.0, 1.0, 51)
\end{lstlisting}
\begin{itemize}
    \item If assets are included:
    \begin{itemize}
        \item loop over \(a'\),
        \item loop over \(L\),
        \item compute \(c\) and leisure \(1-L\),
        \item keep only feasible points.
    \end{itemize}
    \item If assets are not included:
    \begin{itemize}
        \item do not force an asset dimension into the model,
        \item solve the simpler lecture version instead.
    \end{itemize}
    \item Save policies for labor, consumption, and assets if applicable.
\end{itemize}

\textbf{Important detail:}
\begin{itemize}
    \item Make sure leisure is coded as \(1-L\) and stays positive.
\end{itemize}

\textbf{Deliverable:}
\begin{itemize}
    \item A labor-supply module that can explain both substitution and wealth-effect intuition.
\end{itemize}

\textbf{Quick check:}
\begin{itemize}
    \item Labor always lies in \([0,1]\), and the resulting plots are economically interpretable.
\end{itemize}

\subsection{Step 12: Unify simulation, forecasting, and moments across all three models}
\textbf{Goal:} Give the app one common interface across all models.

\textbf{Do this:}
\begin{itemize}
    \item Make sure each model returns results in a comparable structure.
    \item Standardize naming for states, controls, and series.
    \item Have one shared simulation interface and one shared forecast interface.
    \item Build a shared moments table generator that can adapt to different variables by model.
    \item Decide which series are the key variables for each model.
\end{itemize}

\textbf{Example workflow:}
\begin{itemize}
    \item Make each model return a dictionary with the same top-level keys:
    \begin{itemize}
        \item \verb|solution|,
        \item \verb|simulation|,
        \item \verb|forecast|,
        \item \verb|moments|,
        \item \verb|plots_data|.
    \end{itemize}
    \item Use model-specific variable names inside a common format.
    \item This lets \verb|app.py| render different models without branching everywhere.
\end{itemize}

\textbf{Deliverable:}
\begin{itemize}
    \item A common API that the Streamlit app can call without model-specific hacks scattered everywhere.
\end{itemize}

\textbf{Quick check:}
\begin{itemize}
    \item Swapping from Model 1 to Model 2 in the app should mostly change the data, not the architecture.
\end{itemize}

\subsection{Step 13: Build the Streamlit dashboard only after one full model pipeline is stable}
\textbf{Goal:} Put the interface on top of working economics rather than unfinished economics.

\textbf{Do this:}
\begin{itemize}
    \item Load CSS from \verb|assets/style.css|.
    \item Add sidebar navigation between the three models.
    \item Add model-specific parameter controls in the sidebar.
    \item Add controls for transition probabilities, simulation horizon, and forecast horizon.
    \item Solve the selected model using cached functions.
    \item Run simulation and forecasting for the selected model.
    \item Display the outputs in a consistent page structure.
\end{itemize}

\textbf{Example workflow:}
\begin{itemize}
    \item Sidebar controls:
    \begin{itemize}
        \item model selector,
        \item parameters,
        \item transition probabilities,
        \item simulation horizon,
        \item forecast horizon.
    \end{itemize}
    \item Main-page order:
    \begin{itemize}
        \item title and model description,
        \item equations used,
        \item policy function plot,
        \item simulated time-series plot,
        \item forecast plot,
        \item comparative-statics graph,
        \item moments table,
        \item static intuition section,
        \item dynamic summary text.
    \end{itemize}
    \item Cache the expensive solver call so minor UI interactions do not trigger full re-solves unnecessarily.
\end{itemize}

\textbf{Deliverable:}
\begin{itemize}
    \item One app that makes the three models feel like one project rather than three unrelated scripts.
\end{itemize}

\textbf{Quick check:}
\begin{itemize}
    \item A user can switch models without the app breaking or looking inconsistent.
\end{itemize}

\subsection{Step 14: Add comparative statics and dynamic written summaries}
\textbf{Goal:} Make the app explain what changes when parameters change.

\textbf{Do this:}
\begin{itemize}
    \item Add a baseline-versus-alternative parameter comparison for at least one model.
    \item Good candidates are higher \verb|beta|, higher \verb|sigma|, or different wage or TFP persistence.
    \item Plot the two policy functions on the same figure or in side-by-side panels.
    \item Write summary functions that read the moments and generate deterministic interpretation text.
\end{itemize}

\textbf{Example workflow:}
\begin{itemize}
    \item Solve one baseline calibration and one alternative calibration.
    \item Example:
    \begin{itemize}
        \item baseline \(\beta = 0.95\),
        \item alternative \(\beta = 0.98\).
    \end{itemize}
    \item Plot both policy functions together and label them clearly.
    \item Generate summary text with templates such as:
\end{itemize}
\begin{lstlisting}
summary = (
    f"Mean consumption is {mean_c:.2f}. "
    f"When beta rises, the policy shifts toward more saving."
)
\end{lstlisting}
\begin{itemize}
    \item Keep the runtime text deterministic and tied to computed results.
\end{itemize}

\textbf{A good summary should do four things:}
\begin{itemize}
    \item identify the variable,
    \item state what changed,
    \item connect it to economic intuition,
    \item avoid claiming more than the computed results show.
\end{itemize}

\textbf{Deliverable:}
\begin{itemize}
    \item Visual and written interpretation of live results.
\end{itemize}

\textbf{Quick check:}
\begin{itemize}
    \item The text below the charts changes when the user changes parameters, and the wording reflects the actual computed moments.
\end{itemize}

\subsection{Step 15: Add styling, caching, and performance protection}
\textbf{Goal:} Make the app pleasant to demo and fast enough to use live.

\textbf{Do this:}
\begin{itemize}
    \item Create a dark or navy background and bright white content cards.
    \item Use a robust serif font stack if local Garamond is unavailable.
    \item Wrap expensive solver functions in Streamlit caching.
    \item Choose default grid sizes and horizons that update within a few seconds.
    \item Keep plot labels clean and readable.
    \item Make sure the page layout is visually consistent across models.
\end{itemize}

\textbf{Example workflow:}
\begin{itemize}
    \item Put visual rules in \verb|assets/style.css|.
    \item Start with conservative defaults such as:
    \begin{itemize}
        \item 150--250 grid points,
        \item 100--200 simulated periods,
        \item cached solver output keyed by parameter values.
    \end{itemize}
    \item Time common interactions such as changing \(\beta\) or forecast length.
    \item If one slider creates a long delay, lower the default grid density or cache more aggressively.
\end{itemize}

\textbf{Deliverable:}
\begin{itemize}
    \item A dashboard that looks intentional and behaves smoothly.
\end{itemize}

\textbf{Quick check:}
\begin{itemize}
    \item Standard slider changes should update quickly on an ordinary laptop.
\end{itemize}

\subsection{Step 16: Test mathematical correctness, numerical stability, and UI behavior}
\textbf{Goal:} Convert ``works on my machine'' into ``safe to demo live.''

\textbf{Do this:}
\begin{itemize}
    \item Add tests to verify positive consumption, valid policy bounds, and convergence.
    \item Check for NaNs, infinities, or jagged policy plots.
    \item Verify that interpolation never pushes states outside valid ranges.
    \item Confirm that comparative statics move in sensible directions.
    \item Test every page, slider, and forecast control in the dashboard.
    \item Compare the moments table to direct calculations from the simulated data.
\end{itemize}

\textbf{Example workflow:}
\begin{itemize}
    \item Add tests such as:
    \begin{itemize}
        \item ``all simulated consumption values are positive,''
        \item ``all labor choices lie in \([0,1]\),''
        \item ``solver converges before max iterations.''
    \end{itemize}
    \item Inspect plots visually for jagged or discontinuous policies.
    \item Compare dashboard moments to direct \verb|numpy| or \verb|pandas| calculations from the same arrays.
    \item Click through all pages and try awkward inputs intentionally.
\end{itemize}

\textbf{Deliverable:}
\begin{itemize}
    \item A tested project rather than a fragile one.
\end{itemize}

\textbf{Quick check:}
\begin{itemize}
    \item Every model passes a basic suite of solver, simulation, and UI sanity checks.
\end{itemize}

\subsection{Step 17: Keep a prompt log and explainable comments throughout development}
\textbf{Goal:} Be ready to explain both the code and the AI workflow.

\textbf{Do this:}
\begin{itemize}
    \item Keep a simple prompt log with one entry per major component.
    \item For each entry, record the prompt, what file it helped create, and how the output was checked or corrected.
    \item Add comments explaining the economics, not just the syntax.
    \item Mark the key lines in the code that implement constraints, Bellman updates, and policy extraction.
\end{itemize}

\textbf{Example workflow:}
\begin{itemize}
    \item A prompt-log entry should contain:
    \begin{itemize}
        \item prompt used,
        \item file affected,
        \item what part of the output was kept,
        \item what had to be corrected,
        \item how the economics were verified.
    \end{itemize}
    \item In code comments, prefer comments like ``Compute expected continuation value over next-period shocks'' instead of comments that only restate syntax.
\end{itemize}

\textbf{Deliverable:}
\begin{itemize}
    \item Evidence that AI was used responsibly and verified.
\end{itemize}

\textbf{Quick check:}
\begin{itemize}
    \item If the professor points to a function, someone in the group can explain both what it does and how it was validated.
\end{itemize}

\subsection{Step 18: Rehearse the live demo exactly as it will happen}
\textbf{Goal:} Practice the actual grading situation, not just the software.

\textbf{Do this:}
\begin{itemize}
    \item Start from a fresh terminal and run the full install and launch sequence.
    \item Practice switching between all three models.
    \item Practice changing parameters and narrating the economic effect.
    \item Practice explaining why VFI is used, what the Markov chain does, why interpolation is needed, and why caching matters.
    \item Practice answering ``show me where this is in the code.''
\end{itemize}

\textbf{Example workflow:}
\begin{itemize}
    \item Rehearsal round 1:
    \begin{itemize}
        \item one teammate runs the terminal,
        \item one teammate acts as the professor,
        \item one teammate records confusion points.
    \end{itemize}
    \item Rehearsal round 2:
    \begin{itemize}
        \item switch presenters,
        \item ask code questions like ``where is the Bellman update?'' and ``where is interpolation used?''
    \end{itemize}
    \item Rehearsal round 3:
    \begin{itemize}
        \item rapidly switch models and change parameters to test robustness.
    \end{itemize}
\end{itemize}

\textbf{Deliverable:}
\begin{itemize}
    \item A group that can both run and defend the project.
\end{itemize}

\textbf{Quick check:}
\begin{itemize}
    \item Any teammate can handle a random question without needing someone else to rescue the explanation.
\end{itemize}

\section{Final Order of Execution}

If the group wants the shortest operational version, follow this exact sequence:

\begin{enumerate}[label=\textbf{\arabic*.}]
    \item Lock equations from the course notes.
    \item Set up the repository and environment.
    \item Create the file skeleton.
    \item Build shared utilities.
    \item Build the generic VFI solver.
    \item Add feasibility checks.
    \item Finish Model 1 completely.
    \item Add simulation and moments for Model 1.
    \item Add forecasting for Model 1.
    \item Implement Model 2 after timing verification.
    \item Implement Model 3 after confirming whether assets are included.
    \item Standardize simulation, forecast, and moments across all models.
    \item Build the Streamlit dashboard.
    \item Add comparative statics and dynamic summaries.
    \item Add styling, caching, and performance tuning.
    \item Run mathematical, numerical, and UI tests.
    \item Maintain the prompt log and comments.
    \item Rehearse the live demo from scratch.
\end{enumerate}

\begin{goalbox}
\textbf{Bottom line.} The project should be built in a way that makes the economics correct first, the numerical methods trustworthy second, and the dashboard polished third. A clean live demo is important, but the strongest defense comes from having a team that can explain exactly why the model equations, policy functions, forecasts, and summaries are correct.
\end{goalbox}

\end{document}
